% ===== 1. はじめに (Introduction) =====
% 役割:
%   - 問題設定: なぜ LLM 応答の意味的整合性評価が必要か
%   - 既存手法の限界 (詳細は related_work へ、ここは概観のみ)
%   - 本研究の貢献を 3 点で箇条書き
%   - 論文構成の予告 (1 段落)
%
% 分量目安: 1 ページ (2 段組換算で約 800 字)
%
% 貢献候補 (locked-in: A 案):
%   PoR/ΔE による LLM 出力評価フレームワークと HA48 における予備検証

\section{はじめに}
\label{sec:introduction}

% TODO: 第一段落 - 問題設定
% LLM 応答の品質評価における既存指標 (BLEU, BERTScore, G-Eval 等) の
% 限界 (構造的整合性と命題カバレッジの分離不能性) を提起

% TODO: 第二段落 - 提案概観
% PoR (Point of Resonance) による 2 軸分離と ΔE 距離指標の導入

% TODO: 第三段落 - 提案概観 + 貢献
% locked-in (Variant 1): 以下の 1 文を第三段落の冒頭に置く
%   「本稿では、LLM 応答の意味的整合性を構造完全性 (S 軸) と命題カバレッジ
%    (C 軸) の 2 軸で分離評価する PoR/$\Delta$E フレームワークを提案し、
%    HA48 における予備検証を通じて人手評価との対応を確認する。」
%
% 主要貢献 (3 項目):
% \begin{enumerate}
%   \item PoR (S 軸 / C 軸) による 2 軸評価フレームワークの提案と、
%         理想応答からの距離指標 $\Delta$E の定式化
%   \item 意味損失関数 $L_{\mathrm{sem}}$ の 3 項分解 (演算子 / 未知性 /
%         因果構造) による判定根拠の解釈可能性
%   \item 48 件の人手アノテーション (HA48) を用いた予備検証による、
%         提案指標と人手評価との対応の実証
% \end{enumerate}
